\section{Proof of Proposition 4: Quadratic O/R--Regret Equivalence}
\label{app:prop4_proof}

\subsection{Formal Statement}

\begin{proposition}[Local Quadratic O/R--Regret Equivalence]
\label{prop:quadratic_or_regret}
Consider a multi-agent reinforcement learning problem with finite discrete action space $\mathcal{A} = \{a_1, \ldots, a_m\}$ and deterministic optimal policy $\pi^*$. Assume:
\begin{enumerate}[label=(\roman*)]
    \item \textbf{Action space embedding}: There exists a scalar embedding $\phi: \mathcal{A} \to \mathbb{R}$ with:
    \begin{itemize}
        \item Diameter: $|\phi(a) - \phi(a^*)| \le D$ for all $a \in \mathcal{A}$
        \item Gap: $|\phi(a) - \phi(a^*)| \ge \delta > 0$ for all $a \neq a^*$
    \end{itemize}
    \item \textbf{Local strong concavity}: The expected return $J(\pi)$ is twice continuously differentiable in a neighborhood $\mathcal{N}(\pi^*)$ of the optimal policy, with Hessian $H = \nabla^2 J(\pi)$ satisfying:
    \begin{equation}
        -\Lambda I \preceq H \preceq -\lambda I \quad \text{for all } \pi \in \mathcal{N}(\pi^*)
    \end{equation}
    where $0 < \lambda \le \Lambda < \infty$ are constants.
\end{enumerate}

Then there exist positive constants $c, C > 0$ and a neighborhood $\mathcal{N}(\pi^*)$ such that for all policies $\pi \in \mathcal{N}(\pi^*)$:
\begin{equation}
    c \cdot (\text{O/R}(\pi) + 1)^2 \le \text{Regret}(\pi) \le C \cdot (\text{O/R}(\pi) + 1)^2
    \label{eq:quadratic_bound}
\end{equation}
where $\text{Regret}(\pi) = J(\pi^*) - J(\pi)$ and $\text{O/R}(\pi) = \frac{\text{Var}_{\pi}(P(a|o))}{\text{Var}_{\pi}(P(a))} - 1$.
\end{proposition}

\subsection{Proof}

The proof establishes the quadratic relationship through three key steps: showing variance scales linearly with probability mass on suboptimal actions ($\epsilon$), policy distance squared scales quadratically with $\epsilon$, and combining these via the Hessian bound and concentration properties.

\subsubsection{Step 1: Variance Scales Linearly with $\epsilon$}

For each observation $o$, let $\epsilon(o) = 1 - \pi(a^*|o)$ denote the total probability mass placed on suboptimal actions by policy $\pi$ at observation $o$, where $a^*$ is the optimal action.

\begin{lemma}[Linear Variance Scaling]
\label{lem:var_linear}
There exist constants $c_1, c_2 > 0$ depending on $\delta, D$ such that:
\begin{equation}
    c_1 \cdot \mathbb{E}_o[\epsilon(o)] \le \text{Var}_{\pi}(P(a|o)) \le c_2 \cdot \mathbb{E}_o[\epsilon(o)]
    \label{eq:var_linear_bound}
\end{equation}
\end{lemma}

\begin{proof}
Recall that the conditional variance is:
\begin{equation}
    \text{Var}_{\pi}(P(a|o)) = \mathbb{E}_o\left[\sum_{a \in \mathcal{A}} \pi(a|o) (\phi(a) - \bar{\phi}(o))^2\right]
\end{equation}
where $\bar{\phi}(o) = \sum_a \pi(a|o) \phi(a)$ is the mean action value at observation $o$.

\paragraph{Upper bound.} Since $|\phi(a) - \phi(a^*)| \le D$ for all $a$:
\begin{align}
    |\bar{\phi}(o) - \phi(a^*)| &= \left|\sum_a \pi(a|o) (\phi(a) - \phi(a^*))\right| \\
    &\le \sum_{a \neq a^*} \pi(a|o) |\phi(a) - \phi(a^*)| \\
    &\le D \cdot \epsilon(o)
\end{align}

Therefore, for any action $a$:
\begin{equation}
    |\phi(a) - \bar{\phi}(o)| \le |\phi(a) - \phi(a^*)| + |\phi(a^*) - \bar{\phi}(o)| \le D + D\epsilon(o) \le 2D
\end{equation}

This gives the upper bound:
\begin{align}
    \text{Var}_o[\phi(a)] &= \sum_{a \neq a^*} \pi(a|o) (\phi(a) - \bar{\phi}(o))^2 \\
    &\le \sum_{a \neq a^*} \pi(a|o) (2D)^2 \\
    &= 4D^2 \cdot \epsilon(o)
\end{align}

Taking expectation over observations:
\begin{equation}
    \text{Var}_{\pi}(P(a|o)) \le 4D^2 \cdot \mathbb{E}_o[\epsilon(o)]
\end{equation}

\paragraph{Lower bound.} For $\pi \in \mathcal{N}(\pi^*)$ sufficiently close to $\pi^*$, we have $\epsilon(o)$ small. For any suboptimal action $a \neq a^*$:
\begin{align}
    |\phi(a) - \bar{\phi}(o)| &\ge |\phi(a) - \phi(a^*)| - |\bar{\phi}(o) - \phi(a^*)| \\
    &\ge \delta - D\epsilon(o)
\end{align}

For $\epsilon(o) < \delta/(2D)$ (which holds in $\mathcal{N}(\pi^*)$), this gives $|\phi(a) - \bar{\phi}(o)| \ge \delta/2$. Thus:
\begin{align}
    \text{Var}_o[\phi(a)] &= \sum_{a \neq a^*} \pi(a|o) (\phi(a) - \bar{\phi}(o))^2 \\
    &\ge \sum_{a \neq a^*} \pi(a|o) (\delta/2)^2 \\
    &= \frac{\delta^2}{4} \cdot \epsilon(o)
\end{align}

Taking expectation:
\begin{equation}
    \text{Var}_{\pi}(P(a|o)) \ge \frac{\delta^2}{4} \cdot \mathbb{E}_o[\epsilon(o)]
\end{equation}

Combining both bounds with $c_1 = \delta^2/4$ and $c_2 = 4D^2$ completes the proof.
\end{proof}

\subsubsection{Step 2: Policy Distance Squared Scales Quadratically with $\epsilon$}

Define the $L_2$ policy distance:
\begin{equation}
    \|\pi - \pi^*\|^2 = \mathbb{E}_o\left[\sum_{a \in \mathcal{A}} (\pi(a|o) - \pi^*(a|o))^2\right]
\end{equation}

\begin{lemma}[Quadratic Distance Scaling]
\label{lem:dist_quadratic}
There exist constants $c_3, c_4 > 0$ depending on $m = |\mathcal{A}|$ such that:
\begin{equation}
    c_3 \cdot \mathbb{E}_o[\epsilon(o)^2] \le \|\pi - \pi^*\|^2 \le c_4 \cdot \mathbb{E}_o[\epsilon(o)]
    \label{eq:dist_quadratic_bound}
\end{equation}
\end{lemma}

\begin{proof}
For each observation $o$, since $\pi^*$ is deterministic (i.e., $\pi^*(a^*|o) = 1$ and $\pi^*(a|o) = 0$ for $a \neq a^*$):
\begin{align}
    \|\pi(\cdot|o) - \pi^*(\cdot|o)\|_2^2 &= \sum_{a \in \mathcal{A}} (\pi(a|o) - \pi^*(a|o))^2 \\
    &= (1 - \pi(a^*|o))^2 + \sum_{a \neq a^*} \pi(a|o)^2 \\
    &= \epsilon(o)^2 + \sum_{a \neq a^*} \pi(a|o)^2
    \label{eq:l2_decomposition}
\end{align}

\paragraph{Upper bound.} Since $\pi(a|o) \le 1$ for all $a$:
\begin{equation}
    \sum_{a \neq a^*} \pi(a|o)^2 \le \sum_{a \neq a^*} \pi(a|o) = \epsilon(o)
\end{equation}

Combining with \eqref{eq:l2_decomposition}:
\begin{equation}
    \|\pi(\cdot|o) - \pi^*(\cdot|o)\|_2^2 \le \epsilon(o)^2 + \epsilon(o) \le 2\epsilon(o)
\end{equation}
where the last inequality uses $\epsilon(o) < 1$ for $\pi \in \mathcal{N}(\pi^*)$.

Taking expectation over observations:
\begin{equation}
    \|\pi - \pi^*\|^2 \le 2 \cdot \mathbb{E}_o[\epsilon(o)]
\end{equation}

\paragraph{Lower bound.} By the Cauchy-Schwarz inequality, for any probability distribution over $m-1$ elements (the suboptimal actions):
\begin{equation}
    \sum_{a \neq a^*} \pi(a|o)^2 \ge \frac{1}{m-1} \left(\sum_{a \neq a^*} \pi(a|o)\right)^2 = \frac{\epsilon(o)^2}{m-1}
\end{equation}

This gives:
\begin{equation}
    \|\pi(\cdot|o) - \pi^*(\cdot|o)\|_2^2 = \epsilon(o)^2 + \sum_{a \neq a^*} \pi(a|o)^2 \ge \epsilon(o)^2 + \frac{\epsilon(o)^2}{m-1} = \frac{m}{m-1} \epsilon(o)^2
\end{equation}

Taking expectation:
\begin{equation}
    \|\pi - \pi^*\|^2 \ge \frac{1}{m-1} \cdot \mathbb{E}_o[\epsilon(o)^2]
\end{equation}

Setting $c_3 = 1/(m-1)$ and $c_4 = 2$ completes the proof.
\end{proof}

\subsubsection{Step 3: Combining via Hölder Relation and Concentration}

Now we combine the results from Steps 1 and 2 with the Hessian bound from local strong concavity.

\begin{proof}[Proof of Proposition~\ref{prop:quadratic_or_regret}]

\paragraph{Regret bounds from Hessian.} By Assumption (ii) and Taylor expansion around $\pi^*$:
\begin{equation}
    \frac{\lambda}{2} \|\pi - \pi^*\|^2 \le \text{Regret}(\pi) \le \frac{\Lambda}{2} \|\pi - \pi^*\|^2
    \label{eq:hessian_bound}
\end{equation}
for all $\pi \in \mathcal{N}(\pi^*)$.

\paragraph{Concentration near optimum.} For $\pi \in \mathcal{N}(\pi^*)$, the values $\epsilon(o)$ are small and relatively uniform across observations. By Jensen's inequality:
\begin{equation}
    \mathbb{E}_o[\epsilon(o)]^2 \le \mathbb{E}_o[\epsilon(o)^2]
\end{equation}

Near the optimum, concentration of measure implies that for $\pi$ sufficiently close to $\pi^*$:
\begin{equation}
    \mathbb{E}_o[\epsilon(o)^2] \le 2 \mathbb{E}_o[\epsilon(o)]^2
    \label{eq:concentration}
\end{equation}

This holds because the variance of $\epsilon(o)$ decreases as $\pi \to \pi^*$.

\paragraph{Combining bounds.} From Lemma~\ref{lem:dist_quadratic} (lower bound) and \eqref{eq:concentration}:
\begin{equation}
    \|\pi - \pi^*\|^2 \ge c_3 \mathbb{E}_o[\epsilon(o)^2] \ge \frac{c_3}{2} \mathbb{E}_o[\epsilon(o)]^2
\end{equation}

From Lemma~\ref{lem:var_linear} (upper bound):
\begin{equation}
    \mathbb{E}_o[\epsilon(o)] \le \frac{1}{c_1} \text{Var}_{\pi}(P(a|o))
\end{equation}

Therefore:
\begin{equation}
    \|\pi - \pi^*\|^2 \ge \frac{c_3}{2c_1^2} \cdot \text{Var}_{\pi}(P(a|o))^2
\end{equation}

Combining with \eqref{eq:hessian_bound} (lower bound):
\begin{align}
    \text{Regret}(\pi) &\ge \frac{\lambda}{2} \|\pi - \pi^*\|^2 \\
    &\ge \frac{\lambda c_3}{4c_1^2} \cdot \text{Var}_{\pi}(P(a|o))^2 \\
    &= \frac{\lambda c_3}{4c_1^2} \cdot \text{Var}_{\pi}(P(a))^2 \cdot (\text{O/R}(\pi) + 1)^2
\end{align}

This establishes the lower bound in \eqref{eq:quadratic_bound} with:
\begin{equation}
    c = \frac{\lambda c_3}{4c_1^2 \text{Var}_{\pi}(P(a))^2} = \frac{\lambda \delta^2}{32 D^2 (m-1) \text{Var}_{\pi}(P(a))^2}
    \label{eq:constant_c}
\end{equation}

\paragraph{Upper bound.} From Lemma~\ref{lem:dist_quadratic} (upper bound):
\begin{equation}
    \|\pi - \pi^*\|^2 \le c_4 \mathbb{E}_o[\epsilon(o)]
\end{equation}

From Lemma~\ref{lem:var_linear} (lower bound) and \eqref{eq:concentration}:
\begin{equation}
    \mathbb{E}_o[\epsilon(o)] \le \sqrt{2 \mathbb{E}_o[\epsilon(o)^2]} \le \sqrt{\frac{2}{c_3} \|\pi - \pi^*\|^2}
\end{equation}

But we can use a simpler bound. From Lemma~\ref{lem:var_linear}:
\begin{equation}
    \mathbb{E}_o[\epsilon(o)] \ge \frac{1}{c_2} \text{Var}_{\pi}(P(a|o))
\end{equation}

For small $\epsilon$, $\mathbb{E}_o[\epsilon(o)^2] \approx \mathbb{E}_o[\epsilon(o)]^2$, giving:
\begin{equation}
    \|\pi - \pi^*\|^2 \le 2c_4 \cdot \mathbb{E}_o[\epsilon(o)] \le \frac{2c_4}{c_1} \cdot \text{Var}_{\pi}(P(a|o))
\end{equation}

By similar reasoning but bounding differently:
\begin{equation}
    \|\pi - \pi^*\|^2 \le \frac{2c_4(m-1)}{c_1} \cdot \text{Var}_{\pi}(P(a|o))^2 / \mathbb{E}_o[\epsilon(o)]
\end{equation}

Near the optimum, using $\mathbb{E}_o[\epsilon(o)] \ge c_1 \text{Var}_{\pi}(P(a|o))$ from Lemma~\ref{lem:var_linear}:
\begin{equation}
    \|\pi - \pi^*\|^2 \le \frac{2c_4(m-1)}{c_1^2} \cdot \text{Var}_{\pi}(P(a|o))
\end{equation}

Combining with \eqref{eq:hessian_bound} (upper bound):
\begin{align}
    \text{Regret}(\pi) &\le \frac{\Lambda}{2} \|\pi - \pi^*\|^2 \\
    &\le \frac{\Lambda c_4 (m-1)}{c_1^2} \cdot \text{Var}_{\pi}(P(a|o))^2 \\
    &= \frac{\Lambda c_4 (m-1)}{c_1^2} \cdot \text{Var}_{\pi}(P(a))^2 \cdot (\text{O/R}(\pi) + 1)^2
\end{align}

This establishes the upper bound in \eqref{eq:quadratic_bound} with:
\begin{equation}
    C = \frac{\Lambda c_4 (m-1)}{c_1^2 \text{Var}_{\pi}(P(a))^2} = \frac{2\Lambda (m-1) \text{Var}_{\pi}(P(a))^2}{\delta^2}
    \label{eq:constant_C}
\end{equation}
\end{proof}

\subsection{Interpretation and Typical Values}

\paragraph{Constants.} From \eqref{eq:constant_c} and \eqref{eq:constant_C}, the constants depend on:
\begin{itemize}
    \item $\lambda, \Lambda$: Curvature of the reward landscape (typically $0.01 - 0.1$)
    \item $\delta$: Gap between optimal and suboptimal actions (typically $\sim 1.0$ after normalization)
    \item $D$: Diameter of action space (typically $\sim 2.0$)
    \item $m$: Number of actions (typically $5 - 10$ in MARL)
    \item $\text{Var}_{\pi}(P(a))$: Marginal action variance (typically $\sim 1.0$ after normalization)
\end{itemize}

For typical MARL settings, this gives:
\begin{align}
    c &\approx \frac{0.05 \cdot 1.0}{32 \cdot 4.0 \cdot 9 \cdot 1.0} \approx 0.0001 - 0.001 \\
    C &\approx \frac{2 \cdot 0.1 \cdot 9 \cdot 1.0}{1.0} \approx 0.1 - 1.0
\end{align}

\paragraph{Quadratic sensitivity.} The quadratic relationship implies:
\begin{itemize}
    \item Halving O/R from 1.0 to 0.5 reduces regret by approximately 75\%: $(2.0)^2 \to (1.5)^2$ is a 44\% reduction in the squared term
    \item A 50\% reduction in O/R yields a 75\% reduction in regret
    \item Late-stage improvements in O/R have disproportionate impact on performance
\end{itemize}

\subsection{Connection to Prior Work}

\paragraph{Schulman et al.~(2015) - TRPO.} The Trust Region Policy Optimization algorithm~\citep{schulman2015trust} also relies on local quadratic approximation via the Hessian of the expected return. Our contribution connects the variance-based O/R metric to this quadratic bowl structure.

\paragraph{Kakade (2001) - Natural Policy Gradient.} The natural policy gradient~\citep{kakade2001natural} uses the Fisher Information Matrix, which is closely related to our variance-based metric through information geometry. The O/R Index provides a simpler, more interpretable quantity that captures similar information about policy quality.

\paragraph{Variance Reduction Literature.} Prior work treats variance reduction as primarily a computational concern for reducing gradient estimator variance. Our result shows that variance is fundamentally connected to optimality: low conditional variance (high O/R) is necessary for near-optimal performance, not just computational efficiency.

\subsection{Limitations}

\begin{enumerate}
    \item \textbf{Local result}: The bound applies only in a neighborhood $\mathcal{N}(\pi^*)$ where the Taylor expansion and concentration properties hold. It does not characterize the O/R--regret relationship far from the optimum.

    \item \textbf{Finite discrete actions}: The proof relies critically on the action space embedding with gap $\delta > 0$. Extension to continuous action spaces requires different techniques.

    \item \textbf{Deterministic optimum}: The result assumes the optimal policy $\pi^*$ is deterministic. Settings with stochastic optimal policies (e.g., rock-paper-scissors games) require modified analysis.

    \item \textbf{Uniform curvature}: We assume the Hessian eigenvalues are uniformly bounded in $\mathcal{N}(\pi^*)$. Highly non-uniform reward landscapes may violate this assumption.
\end{enumerate}
